% Section 4 (group B): particles from C (x) H (x) O. Sources: Companion thm:SM, thm:octonion-SU3, thm:three-generations,
% rounds v4.9.0-r34, r35, r37, r40--r43, r48, r49; v4.10.0-r13, r14.4, r18.4, r19, r25--r27, r59, r60.
\section{Particles from \texorpdfstring{$\mathbb C\otimes\mathbb H\otimes\mathbb O$}{C x H x O}: a conditional construction}
\label{sec:B-particles}

The carrier of Section~\ref{sec:A-carrier} has a natural extension to the Dixon algebra
$W=\mathbb C\otimes\mathbb H\otimes\mathbb O\cong\mathbb C^{32}$ \cite{Dixon1994}. The Companion states in its main text that
this extension ``generates'' the Standard Model gauge group (\cite{QSC}, \texttt{thm:SM}) and that the three generations
sit on three octonion units (\cite{QSC}, \texttt{thm:three-generations}). Both statements carry status notes in the
Companion today. The first is too strong: the complex scalar alone generates a global phase, not electromagnetism. The
second places colour and generation on the same object, because the cyclic map $e_4\to e_5\to e_6$ is a colour rotation.
This section reports what survives the later rounds. The result is \Conditional{} throughout: once a small number of named
choices is made, one generation and its gauge quantum numbers follow, while the number of generations, the masses and the
mixing do not.

\subsection{One generation under named choices}
\label{sec:B-one-generation}

\paragraph{Colour and charge.}
Let $L_a$ and $R_a$ denote left and right multiplication by $e_a$ on $\mathbb O$, and let $i$ be the $\mathbb C$ factor of $W$,
read as the quantum-mechanical unit (the ``outer'' unit, as opposed to the ``inner'' unit $L_{e_7}$). Colour is the
stabiliser $SU(3)_c=\operatorname{Stab}_{G_2}(e_7)$ of one imaginary unit (\cite{QSC}, \texttt{thm:octonion-SU3}). Following
Furey \cite{Furey2015}, put
\begin{equation}
 \alpha_a=\tfrac12\bigl(L_{e_a}+i\,L_{e_7e_a}\bigr)\quad(a=4,5,6),\qquad
 N=\sum_a\alpha_a^\dagger\alpha_a,\qquad Q=\tfrac13N .
 \label{eq:B-number-operator}
\end{equation}
The Clifford relations of the $L_a$ give $\alpha_a^2=0$ and $\{\alpha_a,\alpha_b^\dagger\}=\delta_{ab}$, so $N$ is a number
operator with eigenvalues $0,1,2,3$ of multiplicities $1,3,3,1$. Every element of $SU(3)_c$ commutes with $Q$; the
eigenspaces $Q=0$ and $Q=1$ are colour singlets spanned by $1\pm ie_7$, the eigenspaces $Q=\tfrac13$ and $Q=\tfrac23$ are
conjugate colour triplets (\rndold{r34}). On $\mathbb C\otimes\mathbb O$ the same sectors are the joint eigenspaces of the
commuting involutions $iL_{e_7}$ and $iR_{e_7}$,
\begin{equation}
 P_{(\ell,r)}=\tfrac14\,(1+\ell\,iL_{e_7})(1+r\,iR_{e_7}),\qquad 6Q=3-iL_{e_7}-2\,iR_{e_7},
 \label{eq:B-sectors}
\end{equation}
with ranks $1,3,3,1$; these tables are derived from the octonion product in Lean~4 (\rndold{r48}). \Established{} as
algebra. It reproduces Furey's charge pattern in the Companion's conventions and is a consistency check, not a prediction.

One property is specific to the frame picture of Section~\ref{sec:A-frames}. The time-reversal part $T$ of the frame maps is complex
conjugation of the $\mathbb C$ factor, and $TQT^{-1}=1-Q$ (\rndold{r34}). The spectrum and the colour structure are
frame-free, but which state carries which charge is not: a charge statement has to name its frame.

\paragraph{Weak isospin, chirality and hypercharge.}
The quaternion factor offers two commuting copies of $\mathfrak{su}(2)$, from the left and from the right. The distribution
used from \rndold{r41} on places the Lorentz group on the left ($L_q$ for rotations, $iL_q$ for boosts) and the weak
$SU(2)$ on the right, $T_a=-\tfrac i2R_{e_a}$. With $Y/2=Q-\tfrac12$ and $Q_{\rm em}=T_3+Y/2$ this gives the hypercharges of
the left-handed doublets and their conjugates, but every state is then a weak doublet and the right-handed singlets are
missing. The missing half is supplied by a projector of the compensator form of Section~\ref{sec:A-carrier} (\rndold{r42}):
\begin{equation}
 C_{\rm ch}=\tfrac12(1+J),\quad J=i\,L_{e_7},\qquad
 T_a'=C_{\rm ch}T_a,\qquad
 \tfrac Y2=Q-\tfrac12+(1-C_{\rm ch})T_3 .
 \label{eq:B-chirality}
\end{equation}
$C_{\rm ch}$ is a Hermitian projector of rank $16$ that commutes with colour, the Lorentz group, the right action of $\mathbb H$
and $Q$. Under $(SU(3)_c,SU(2)',Y/2)$ the space splits into
\begin{equation}
 (3,2,\tfrac16)\oplus(1,2,-\tfrac12)\oplus(\bar3,1,-\tfrac23)\oplus(\bar3,1,\tfrac13)\oplus(1,1,1)\oplus(1,1,0),
 \label{eq:B-content}
\end{equation}
that is $Q_L,L_L,u^c,d^c,e^c,\nu^c$ with complex dimensions $12,4,6,6,2,2$: sixteen Weyl states with two Lorentz
components each, and $Q_{\rm em}=T_3'+Y/2$. The weak $SU(2)'$ acts only on the image of $C_{\rm ch}$; in the model this is
chirality. All four anomaly traces ($\operatorname{Tr}Y$, $\operatorname{Tr}Y^3$, $SU(2)^2U(1)$, $SU(3)^2U(1)$) vanish.

\paragraph{Which choices are forced.}
Round \rndold{r43} asks why this projector and this offset. Fix the setting: $C=\tfrac12(1+J)$ with $J$ Hermitian,
$J^2=1$, $J\neq\pm1$, acting on $\mathbb C\otimes\mathbb O$, commuting with colour and $Q$, of rank 16, with image one
singlet and one triplet; $Q_{\rm em}=T_3+Q+s$. Then:

\begin{proposition}[\rndold{r43}]\label{prop:B-anomaly-choice}
Under this setting there are exactly four chirality projectors, $J=\pm iL_{e_7}$ and $J=\pm iR_{e_7}$. For each,
$\operatorname{Tr}Y=0$ holds exactly for $s=-\tfrac12$. At $s=-\tfrac12$ all four anomaly traces vanish exactly for
$J=\pm iL_{e_7}$, while $J=\pm iR_{e_7}$ give $\operatorname{Tr}Y^3=\pm\tfrac32$. The two admissible choices are frame images
of each other, $T(iL_{e_7})T^{-1}=-iL_{e_7}$.
\end{proposition}

\begin{proof}[Proof sketch]
The commutant of colour and $Q$ on $\mathbb C\otimes\mathbb O$ consists of scalars on the four $Q$-sectors, since each is
colour-irreducible with its own charge. A Hermitian involution there is a sign per sector, and the rank and image conditions
leave four. The traces are exact fractions over the sixteen states (script \texttt{k84}).
\end{proof}

Gravitational anomaly cancellation fixes the hypercharge offset, and the remaining anomalies select left against right
multiplication. What is left is the Standard Model and its frame mirror. The hypercharge offset also has a clean reading
(\rndold{r49}): on one generation $2(Q-\tfrac12)=B-L$, hence
\begin{equation}
 \tfrac Y2=\tfrac{B-L}{2}+T_3^R,\qquad Q_{\rm em}=T_3^L+T_3^R+\tfrac{B-L}{2},\qquad T_3^L=C\,T_3,\ T_3^R=(1-C)\,T_3 ,
 \label{eq:B-BL}
\end{equation}
the left-right symmetric charge formula. The axis $e_7$ carries $B-L$, not the electric charge: no combination of
$1$, $iL_{e_7}$, $iR_{e_7}$ and their product acts as $Q_{\rm em}$ on a generation (it misses by the $T_3$ splitting inside a
doublet). The photon needs the weak isospin from $\mathbb H$ as well.

Table~\ref{tab:B-choices} separates what is set from what follows. None of this is a prediction: that the Standard Model
content is anomaly free is known, and similar constructions on $\mathbb C\otimes\mathbb H\otimes\mathbb O$ exist in the
literature \cite{Dixon1994,Furey2015}. What the Companion adds is the placement on the frames and sides of its own carrier,
and the observation that, in this structure, anomaly cancellation alone decides between left and right.

\begin{table}[t]
\centering\small
\begin{tabularx}{\textwidth}{@{}>{\raggedright\arraybackslash}X>{\raggedright\arraybackslash}X@{}}
\toprule
set (choices) & follows (under the choices)\\
\midrule
$\mathbb C$ factor as the quantum-mechanical unit; colour as $\operatorname{Stab}_{G_2}(e_7)$ &
charge spectrum $0,\tfrac13,\tfrac23,1$ with multiplicities $1,3,3,1$; singlets and conjugate triplets (\rndold{r34}, \rndold{r48})\\
Lorentz group on the left of $\mathbb H$, weak $SU(2)$ on the right &
doublet structure; $T$ acts as charge conjugation $Q_{\rm em}\mapsto-Q_{\rm em}$ (\rndold{r41})\\
chirality projector of compensator form on $\mathbb C\otimes\mathbb O$; $Q_{\rm em}=T_3+Q+s$ &
$J=\pm iL_{e_7}$ and $s=-\tfrac12$ from anomaly cancellation; sixteen Weyl states \eqref{eq:B-content} (\rndold{r42}, \rndold{r43})\\
--- & $e_7$ carries $B-L$; charge formula \eqref{eq:B-BL} (\rndold{r49})\\
number of generations; symmetry breaking; couplings & not addressed\\
\bottomrule
\end{tabularx}
\caption{One generation on $\mathbb C\otimes\mathbb H\otimes\mathbb O$: choices and consequences. Status \Conditional.}
\label{tab:B-choices}
\end{table}

\paragraph{No ban on a fourth generation.}
The search for a second threefold structure outside $\mathbb O$ led to triality (\rndold{r37}): the three slots of
$\mathbb O^3$ carry the three inequivalent eight-dimensional representations of $\mathrm{Spin}(8)$, the cyclic slot shift is
the outer triality automorphism, and what it fixes is $G_2\supset SU(3)_c$. Generation would then be what the traversal
permutes and colour what it fixes. That there are exactly three slots rests on $\operatorname{Out}(\mathrm{Spin}(8))=S_3$.
This excludes a fourth slot with a new eight-dimensional representation, not a repetition of one of the three. The
structure therefore does not forbid a fourth generation, and that the slots \emph{are} the generations is a \Reading.

\subsection{A generation qubit needs copies}
\label{sec:B-commutant}

Before asking where three generations could sit, one can ask which operators on one generation could distinguish
generations at all. They must commute with the gauge and Lorentz generators.

\begin{proposition}[Commutant of one generation, \rnd{r19}]\label{prop:B-commutant}
On $W=\mathbb C\otimes\mathbb H\otimes\mathbb O$ with the Lorentz generators $L_q$, $iL_q$, the weak generators $T_a'$, the
hypercharge of \eqref{eq:B-chirality} and the eight colour derivations, the complex-linear operators commuting with all
generators form a commutative algebra $\cong\mathbb C^6$, one phase per multiplet of \eqref{eq:B-content}. No antilinear
operator commutes with all generators. On $n$ copies of the generation the commutant has dimension $6n^2$ and contains
$M_n(\mathbb C)$ acting on the copy index.
\end{proposition}

\begin{proof}
The six multiplets in \eqref{eq:B-content} are irreducible, pairwise inequivalent and of complex type, each with
multiplicity one. By Schur's lemma a commuting complex-linear operator is a scalar on each. On $n$ copies each multiplet
$\mathcal K_a$ appears as $\mathcal K_a\otimes\mathbb C^n$, and the commutant is $\bigoplus_a\mathrm{id}_{\mathcal K_a}\otimes M_n(\mathbb C)$.
The explicit computation (\rnd{r19}.1, \rnd{r25}.1) confirms this on the concrete carrier, including the absence of an
antilinear part.
\end{proof}

The consequence is \Established: on one generation there is no pair of anticommuting symmetric operators, hence no qubit
with three mutually unbiased bases, whether built from the frame maps or from anything else. Of the frame maps,
$R\colon z\mapsto iz$ is central and $T$, $J_B$ map a generation to its conjugate. A generation index exists only when the
carrier contains the generation several times, and the number of copies is then an input (\rnd{r19}.2). The reading ``three
generations are the three frame maps'' of Section~\ref{sec:B-koide} therefore takes the number three as input on any
charged carrier.

\subsection{Three routes to a generation index}
\label{sec:B-routes}

Three carriers were examined for a neutral generation structure with an order-three cycle and a reference state at the
``magic'' angle $\arccos(1/\sqrt3)$, the profile that would produce the Koide octant of Section~\ref{sec:B-koide}.
Figure~\ref{fig:B-routes} summarises what each delivers. All items are explorations, \Exploration, except where a test
is named.

\paragraph{Route A: sedenions with charge axis $e_8$.}
On $\mathbb C\otimes\mathbb S$ (sedenions $\mathbb S$, cycle $\psi$ of order three) a backward search over all charges with
$[Q,\psi]=0$, $[Q,\mathfrak{su}(3)_c]=0$, $\bar Q=1-Q$ and spectrum $\{0,\tfrac13,\tfrac23,1\}$ gives, up to two signs,
\begin{equation}
 Q_*=P_0\bigl(\tfrac12-\tfrac i2L_{e_8}\bigr)P_0+\tfrac23P_3+\tfrac13P_{\bar3}
 \label{eq:B-Qstar}
\end{equation}
with charge axis $e_8$, the only $\psi$-fixed imaginary direction (\rnd{r13}.1). The axis is an assumption, registered
as placeholder CT-003; everything below inherits it. The kernel $N=\operatorname{span}(1+ie_8,\,e_7-ie_{15})$ is a colour-
and charge-neutral qubit, on which $\psi$ acts as $\operatorname{diag}(1,\omega)$. The quaternion algebra
$\mathbb H_n=\operatorname{span}(1,e_7,e_8,e_{15})$ acts on $N$ from the right; its 24 Hurwitz units act as the binary
tetrahedral group $2T$, whose four threefold axes lie at $|\cos|=1/\sqrt3$ to the $\psi$ axis. An eigenstate of an
order-three Hurwitz unit has orbit overlap $\rho^2=\tfrac12$, hence $K=\tfrac23$ and Bargmann phase $\pi/4$ (\rnd{r13}.2).
The target was known in advance: this is a construction, not a derivation. The decisive negative is the content: under
$(Q_*,\text{colour},\psi)$ the space carries exactly two complete generations, and $\psi$ acts inside each as a phase, not as a
generation step (\rnd{r13}.3). The Majorana vector on $N$ is not fixed, and two flavours do not give three masses
(\rnd{r13}.5). A test of the inherited charge (H103a) was \Refuted.

\paragraph{Route B: triality slots.}
On $(\mathbb O^3,\tau)$ with $\tau(x,y,z)=(y,z,x)$ there are three slots and a cycle. With the slot charge every species
splits $\tau$-symmetrically as $2+1$, and on the doublet $\tau=\operatorname{diag}(\omega,\bar\omega)$. But a real slot vector
lies on the equator of that doublet and gives $\rho^2=\tfrac14$, $K=\tfrac12$; and no operation at hand supplies a twofold
axis at the magic angle (of $10\,368$ monomial slot operations with twelfth roots of unity, $72$ preserve the splitting, and
they rotate only about the $\tau$ axis or perpendicular to it) (\rnd{r13}.6). The two routes are complementary: the
quaternion group needs a doublet made of two slots (route A), the slot cycle needs three slots and turns the doublet into a
Fourier pair (route B). With the present inventory the profile cannot be met.

\paragraph{Dixon copies on the Hilbert-shift tower.}
Proposition~\ref{prop:B-commutant} suggests copies. Read as stages of the Companion's Hilbert-shift tower, two Dixon
algebras form a stage HS(7) and four form HS(8) (\rnd{r25}). The duplication rule above HS(6) is a setting; under it both
stages are admissible, and all copy operations are free of leakage and debt (\rnd{r26}, with the correction \rnd{r27}).
Two copies give a copy qubit, whose three unbiased states form the flat octant of Section~\ref{sec:B-koide}, with one
massless generation. Four copies give $\mathbb C^4=\mathrm{Sym}^2\oplus\Lambda^2$, a triplet and a singlet; the spin-one
triangle of Section~\ref{sec:B-koide} lies in the triplet, and the only equivariant map from three triplet vectors to the
singlet is the triple product $\det V$, with $|\det V|^2=\det G$. The singlet is occupied exactly when all three masses are
positive. This is a \Reading: the number of copies is set, and that the generations are read out as states with overlap
$\tfrac12$ remains the Koide input.

% F4: what the generation carriers deliver (Companion v4.10.0-r13, r19, r25, r26, r59, r60; v4.9.0-r37).
% Cell colours: status colours of the article; grey = not examined / not supplied.
\begin{figure}[t]
\centering
\begin{tikzpicture}[x=1cm,y=1cm,
  cell/.style={draw=white,line width=1.2pt,minimum width=3.2cm,minimum height=1.5cm,text width=3.0cm,
               align=center,font=\scriptsize,inner sep=2pt,anchor=north west},
  head/.style={cell,font=\scriptsize\bfseries,fill=black!8},
  rowh/.style={cell,minimum width=2.6cm,text width=2.4cm,align=left,font=\scriptsize\bfseries,fill=black!8},
  est/.style={cell,fill=st-est!22},
  cond/.style={cell,fill=st-cond!22},
  read/.style={cell,fill=st-read!20},
  ref/.style={cell,fill=st-ref!20},
  open/.style={cell,fill=st-open!22},
  na/.style={cell,fill=black!4,text=black!60}]
\def\cw{3.2}\def\rh{1.5}\def\hx{2.6}\def\hh{1.25}
\newcommand{\C}[4]{\node[#3] at ({\hx+(#1)*\cw},{-\hh-(#2-1)*\rh}) {#4};}
\newcommand{\Hd}[2]{\node[head,minimum height=\hh cm] at ({\hx+(#1)*\cw},0) {#2};}
\newcommand{\Rw}[2]{\node[rowh] at (0,{-\hh-(#1-1)*\rh}) {#2};}
% header row
\node[rowh,minimum height=\hh cm] at (0,0) {carrier};
\Hd{0}{Route A\\ sedenions $\mathbb C\otimes\mathbb S$,\\ axis $e_8$ (CT-003)}
\Hd{1}{Route B\\ triality slots $\mathbb O^3$,\\ cycle $\tau$}
\Hd{2}{Dixon copies\\ HS(7), HS(8)\\ ($n=2,4$ copies, set)}
\Hd{3}{Albert algebra\\ $H_3(\mathbb O)$, slots\\ $x_1,x_2,x_3$}

% row 1: family symmetry
\Rw{1}{family symmetry}
\C{0}{1}{cond}{$\mathfrak{su}(2)_\ell\oplus\mathfrak{su}(2)_q$ from the right action of $\mathbb H_n$ (\rnd{r13}.4)}
\C{1}{1}{est}{$S_3$ on the slots (\rndold{r37})}
\C{2}{1}{est}{$M_n(\mathbb C)$ on the copy index (\rnd{r19}, \rnd{r25})}
\C{3}{1}{est}{$SO(3)$ derived; charge breaks it to exactly $A_4$, pure spin 3 (\rnd{r60})}
% row 2: three generations
\Rw{2}{three generations?}
\C{0}{2}{ref}{no: exactly two (\rnd{r13}.3)}
\C{1}{2}{read}{three slots, read as generations; fourth not forbidden (\rndold{r37})}
\C{2}{2}{cond}{only as many as copies are set (\rnd{r19}.2)}
\C{3}{2}{read}{three slots, read as generations (\rnd{r60})}
% row 3: cycle on a neutral qubit
\Rw{3}{order-3 cycle on a neutral qubit?}
\C{0}{3}{cond}{$\psi=\mathrm{diag}(1,\omega)$ on $N$, but a phase inside the generations (\rnd{r13}.1, .3)}
\C{1}{3}{est}{$\tau=\mathrm{diag}(\omega,\bar\omega)$ on the slot doublet (\rnd{r13}.6)}
\C{2}{3}{na}{not supplied: frame maps do not act on copies (\rnd{r25})}
\C{3}{3}{na}{not examined}
% row 4: magic angle
\Rw{4}{reference at the magic angle?}
\C{0}{4}{cond}{yes: Hurwitz units give $2T$, axes at $|\!\cos|=1/\sqrt3$; a construction (\rnd{r13}.2)}
\C{1}{4}{ref}{no operation at hand supplies it (\rnd{r13}.6)}
\C{2}{4}{na}{not examined}
\C{3}{4}{na}{not examined}
% row 5: K = 2/3
\Rw{5}{$K=2/3$?}
\C{0}{5}{cond}{orbit overlap $\tfrac12$, phase $\pi/4$; target known in advance (\rnd{r13}.2)}
\C{1}{5}{ref}{real slot vector: $\rho^2=\tfrac14$, $K=\tfrac12$ (\rnd{r13}.6)}
\C{2}{5}{cond}{for chosen states only: flat octant ($n=2$), spin-1 triangle ($n=4$)}
\C{3}{5}{na}{not addressed}
% row 6: mixing
\Rw{6}{mixing?}
\C{0}{6}{open}{Majorana vector on $N$ not fixed; no three masses (\rnd{r13}.5)}
\C{1}{6}{open}{TM2 under an assumed $\mathbb Z_2$; $\theta_{23}$ free (\rnd{r18}.4)}
\C{2}{6}{na}{not addressed}
\C{3}{6}{open}{exact $A_4$: $V_{\rm CKM}=I$; alignment depends on coefficients (\rnd{r60})}
% legend
\begin{scope}[yshift=-10.7cm,font=\scriptsize]
\foreach \c/\t/\x in {st-est/{exact or derived}/0,st-cond/{setting or construction}/2.9,st-read/reading/6.4,st-ref/no/8.2,st-open/open/9.5,black/{not examined}/11.0}{
  \fill[\c!22] (\x,0) rectangle ++(0.35,0.25);
  \node[anchor=west] at (\x+0.4,0.125) {\t};}
\end{scope}
\end{tikzpicture}
\caption{Four candidate carriers for a generation index, and what each delivers. Cell colours give the status in the
Companion: green exact or derived, yellow under a named setting or a construction whose target was known, purple a
reading, red a negative answer, orange open, grey not examined. No column delivers three generations, an order-three cycle,
the magic-angle reference, $K=2/3$ and a mixing pattern at once. Route~A and route~B are complementary
(\rnd{r13}.6): the quaternion group needs a doublet made of two slots, the slot cycle needs three slots.}
\label{fig:B-routes}
\end{figure}


\subsection{The Albert algebra and \texorpdfstring{$A_4$}{A4}}
\label{sec:B-albert}

Round \rnd{r60} takes the carrier $H_3(\mathbb O)$ of Hermitian $3\times3$ octonion matrices with the Jordan product
$X\circ Y=\tfrac12(XY+YX)$. Its three off-diagonal octonions are the slots, $\operatorname{Aut}=F_4\supset\mathrm{Spin}(8)$
carries triality, colour is the diagonal $SU(3)_c$, and charge is the slot-wise $C=L_{e_7}$. Reading the slots as
generations is a reading; the following statements are algebra (Wolfram exact, scripts \texttt{k275}--\texttt{k277}).

\begin{proposition}[\rnd{r60}.1, .3, .4]\label{prop:B-albert}
\begin{enumerate}
\item[(a)] The commutant of $G_2$ in $F_4$ is the real $SO(3)$ acting by $X\mapsto RXR^{\mathsf T}$. Under it every colour
triplet is spin $1$; the colour singlets decompose as $1\oplus3\oplus5$.
\item[(b)] Charge breaks this family $SO(3)$ to exactly $A_4$, with the slots in its three-dimensional irreducible
representation.
\item[(c)] On the colour triplets every $A_4$-invariant operator is $SO(3)$-invariant. Hence, with exact $A_4$, quark masses
are degenerate and $V_{\rm CKM}=I$.
\item[(d)] For a flavon $v$ in the $e_7$ parts of the slots, $Y=C\,X(v)$ is real symmetric with zero diagonal and entries
$-v_k$, and its cubic Albert norm is $N(Y)=\det Y=-2\,v_1v_2v_3$, with $\operatorname{Tr}Y=0$, $\operatorname{Tr}Y^2=2|v|^2$.
\end{enumerate}
\end{proposition}

\begin{proof}[Proof sketch]
(b) Of the 24 signed permutations of determinant $+1$, exactly the 12 with even slot permutation commute with charge; the
odd ones exchange $x_{ij}$ with $\bar x_{ij}$, and conjugation turns $L_{e_7}$ into $-R_{e_7}$. The residual Lie algebra is
zero, and every finite subgroup of $SO(3)$ containing $A_4$ is $A_4$, $S_4$ or $A_5$; an icosahedral rotation breaks.
(c) Spin $1\otimes1$ contains $l=0,1,2$ only, and $A_4$ has no invariant at $l=1$ or $l=2$. Since $C$ is $A_4$-invariant, any
$SO(3)$-equivariant construction from $C$ stays $SO(3)$-invariant on the triplets. (d) A real symmetric matrix with zero
diagonal and off-diagonal entries $a,b,c$ has determinant $2abc$.
\end{proof}

The $A_4$ that earlier rounds set by hand as a slot symmetry is here the remnant of a derived family symmetry. The test
H127 asked which spin carries the breaking. Stage~A was \textbf{confirmed}: the charge part
$A_0=P_sCP_s$ lies entirely in the spin-3 component ($w_3=1$, $w_{l\neq3}=0$), a single tensor of type $xyz$, the object that
$A_4$ flavon models introduce by hand. Stage~B (a leak into the quarks at some order) was \Refuted{} by its wording, and by
(c) it is empty for every $SO(3)$-compatible coupling. Quark mixing therefore requires that $A_4$ itself breaks.

The vacuum alignment is not fixed. A potential built only from $Y$ depends on $|v|^2$ and $v_1v_2v_3$, and its minimum is of
type $(1,1,1)$ with residual $\mathbb Z_3$; with the mixed quartic invariants the potential is the general $A_4$ potential,
and the minimum depends on the coefficient of $\sum_iv_i^4$: the two sample values computed give type $(1,1,1)$ and type $(1,0,0)$ (residual $\mathbb Z_2$). This is
the known $A_4$ phase diagram, and the coefficients are free. No blind test of the alignment is possible at present
(\rnd{r60}.5). Set remain the slots as generations, the flavon placement, the coefficients, every mass and every scale.

\subsection{A relative CP sign between leptons and quarks}
\label{sec:B-cp}

On route A the right multiplications obey the quaternion relations with opposite orientation on the neutral and on the
coloured part (\rnd{r13}.4), a consequence of the non-associativity of the sedenion product. In terms of the neutral family
$\mathfrak{su}(2)$ (\rnd{r14}.4, under CT-003):
\begin{equation}
 [R_{e_7},R_{e_8}]=-2R_{e_{15}}\ \text{on the lepton block},\qquad [R_{e_7},R_{e_8}]=+2R_{e_{15}}\ \text{on the coloured block}.
 \label{eq:B-cp}
\end{equation}
For three projectors onto the $+1$ eigenstates of generators with $[J_x,J_y]=isJ_z$, $s=\pm1$, the Bargmann invariant is
$\operatorname{tr}(P_xP_yP_z)=\tfrac14+\tfrac{is}{4}$. Relative to the Lorentz spin (left action of $\mathbb H$, invariant
$\tfrac14+\tfrac i4$) the lepton doublet is the complex-conjugate one: $\tfrac14-\tfrac i4$ on leptons against
$\tfrac14+\tfrac i4$ on quarks, in all $24+24$ assignments of full rank. This is H125 stage~A, \textbf{confirmed} but
\emph{not blind} (\rnd{r59}.6). The relative sign survives $i\mapsto-i$, which conjugates both sectors at once; it is the
first sign in this series that is not a convention.

Stage~B is a prediction against future data. Under the setting S1 (CP arises in both sectors in the same way from the
family structure) one has $\operatorname{sign}J_{\rm PMNS}=-\operatorname{sign}J_{\rm CKM}$, that is
$\sin\delta_{\rm PMNS}<0$. It is decided by the first analysis that fixes the sign at $\ge3\sigma$, with deadline
31~December~2035; the prior probability is $\tfrac12$, and a confirmation would mainly support S1. \Open.

The absolute handedness is different. Every attempt to fix it in \rnd{r59}.1--.4 set it somewhere, and the open choice
ended at the choice of $i$, which over $\mathbb C$ is a convention (\rnd{r59}.4, \rnd{r59}.15). The orientation of the
Koide octant in Section~\ref{sec:B-koide-orientation} is of the same kind.

\paragraph{Summary of the section.}
Under named choices one generation with its charges, chirality and anomaly freedom follows from $\mathbb C\otimes\mathbb H
\otimes\mathbb O$; this is a consistency result in a known line of work. The number of generations is not derived: one
generation admits no generation qubit, route~A gives two generations, route~B gives three slots without the Koide
reference, copies give whatever number is set, and the Albert algebra gives a derived family $SO(3)$ broken to $A_4$ with
the slots read as generations. Masses, mixing angles and scales are not derived anywhere in this section. OP3 and OP4 remain
\Open.
